\chapter{Algoritmos de trading: ordens, custos e execução}
\label{cap:trading}

\abertura{Um backtest termina em uma curva de capital. Entre essa curva e dinheiro na conta existe uma camada que os capítulos anteriores resumiram em um único número, o custo por ponta: a ordem precisa chegar à bolsa, encontrar contraparte, ser executada a algum preço e pagar por isso. Este capítulo trata dessa camada, do livro de ofertas ao programa que conversa com a corretora, e mostra quanto dela uma estratégia aguenta.}

\begin{objetivos}
  \item descrever os componentes de um bot de trading e o laço de eventos que os liga;
  \item ler um livro de ofertas e escolher entre ordem a mercado, limitada e stop sabendo o que cada uma garante;
  \item estimar o custo de uma operação somando taxas, meio spread e impacto de mercado;
  \item calcular o custo de equilíbrio de uma estratégia e relacioná-lo ao giro;
  \item construir e testar uma corretora simulada sobre barras diárias, conhecendo seus limites;
  \item dimensionar posições e cercar um sistema automático com controles de risco operacional.
\end{objetivos}

\section{Do backtest ao mercado}

Um backtest como os do Capítulo~\ref{cap:backtest} usa a convenção de execução no fechamento do sinal e cobra um custo fixo de 0,10\% por ponta. Por trás disso há três suposições: a ordem é executada no instante em que o sinal fica pronto, pelo preço de fechamento que o gerou, na quantidade inteira. No mercado nenhuma das três vale por construção, e o custo de 0,10\% é uma estimativa que este capítulo ensina a abrir.

\begin{itemize}
  \item \textbf{Latência.} Entre o dado chegar, o programa decidir e a ordem alcançar a bolsa passa um tempo, de milissegundos a minutos conforme a infraestrutura. Nesse intervalo o preço anda.
  \item \textbf{Preço de execução.} Você compra pelo preço de quem aceita vender, que é maior que o de quem aceita comprar, e uma ordem grande consome várias ofertas a preços cada vez piores. A diferença entre o preço que a regra supôs e o preço obtido chama-se \termo{slippage}, ou deslizamento.
  \item \textbf{Preenchimento parcial ou nenhum.} Uma ordem com limite de preço pode ser executada em parte, ou ficar sem execução enquanto o mercado se afasta.
  \item \textbf{Dados sujos.} O histórico que você baixa foi corrigido depois; o fluxo em tempo real chega com negócios cancelados, cotações repetidas e buracos. Nem o histórico chega limpo: a cópia bruta das barras de PETR4, em \arq{dados/bruto/}, tem 54 barras com volume zero e os quatro preços idênticos; o Capítulo~\ref{cap:dados} as identifica e remove as 52 que repetem o fechamento anterior.
  \item \textbf{Falhas.} A conexão cai, a corretora rejeita a ordem, o programa trava com uma posição aberta. Nada disso acontece em um backtest.
\end{itemize}

\citet{chan2021} trata a distância entre resultado simulado e resultado real como parte normal do trabalho e recomenda um período de operação simulada antes de arriscar capital.

\section{A anatomia de um bot}

Um \ingles{bot} de trading é um programa que repete um ciclo: recebe um evento (uma barra nova, um negócio, a confirmação de uma ordem), atualiza o que sabe e decide se envia, altera ou cancela ordens. A Figura~\ref{fig:c11-laco} separa esse ciclo em sete responsabilidades. Em um sistema pequeno elas podem ser sete funções no mesmo arquivo, desde que cada uma faça uma coisa só.

\begin{figure}[t]
\centering
\begin{tikzpicture}[x=1mm, y=1mm, font=\sffamily\footnotesize,
  cx/.style={draw, line width=0.4pt, fill=lilas, minimum width=25mm, minimum height=9.5mm, align=center, inner sep=1pt},
  seta/.style={-{Stealth[length=1.8mm]}, line width=0.4pt}]
  \node[cx] (dados) at (0,0) {Dados};
  \node[cx] (sinal) at (31,0) {Sinal};
  \node[cx] (alvo) at (62,0) {Alvo de\\posição};
  \node[cx] (risco) at (93,0) {Controle\\de risco};
  \node[cx] (ger) at (93,-19) {Gerador\\de ordens};
  \node[cx] (exec) at (62,-19) {Execução};
  \node[cx] (rec) at (31,-19) {Reconciliação\\e registro};
  \node[cx, fill=white, draw=roxo, line width=0.8pt] (cor) at (62,-40) {Corretora\\e bolsa};
  \draw[seta] (dados) -- (sinal);
  \draw[seta] (sinal) -- (alvo);
  \draw[seta] (alvo) -- (risco);
  \draw[seta] (risco) -- (ger);
  \draw[seta] (ger) -- (exec);
  \draw[seta] (exec) -- (rec);
  \draw[seta] (rec.west) -| (dados.south);
  \draw[seta] ([xshift=-3mm]exec.south) -- node[left, text=cinza] {ordens} ([xshift=-3mm]cor.north);
  \draw[seta] ([xshift=3mm]cor.north) -- node[right, text=cinza] {execuções} ([xshift=3mm]exec.south);
  \draw[seta] (cor.west) -| node[pos=0.25, below, text=cinza] {cotações e negócios} ([xshift=-6mm]dados.south);
  \node[text=cinza, anchor=north, align=center] at (9,-19.8) {próximo\\evento};
\end{tikzpicture}
\caption{O laço de eventos de um bot. A linha de cima decide o que se quer ter; a de baixo transforma a decisão em ordens e confere o que de fato aconteceu. A estratégia em si ocupa só a caixa \emph{Sinal}.}
\label{fig:c11-laco}
\end{figure}

A camada de \emph{dados} recebe cotações e negócios, descarta o que é inválido e monta as barras e os indicadores. O \emph{sinal} é a regra da estratégia, a mesma função testada no backtest, e devolve uma opinião: ``comprado'', ou uma previsão de retorno. O \emph{alvo de posição} converte a opinião em quantidade de ações (Seção~\ref{sec:c11-dimensionamento}). O \emph{controle de risco} recusa ou reduz alvos que violem limites (Seção~\ref{sec:c11-risco}). O \emph{gerador de ordens} compara o alvo com a posição atual e com as ordens já abertas e produz a diferença. A \emph{execução} envia as ordens, acompanha seus estados e trata rejeições e execuções parciais. A \emph{reconciliação} confere a posição que o programa julga ter com a que a corretora informa, e o \emph{registro} grava tudo o que aconteceu.

Trabalhar com \emph{alvo de posição}, e não com ``sinal de compra'', evita uma família inteira de erros. Se o programa reinicia no meio do dia, ele relê a posição na corretora, recalcula o alvo e envia só a diferença. Um bot que reage a sinais de compra, ao reiniciar, compra de novo.

\section{Microestrutura mínima: o livro de ofertas}

Nas bolsas eletrônicas, a B3 entre elas, o preço se forma em um \termo{livro de ofertas} (\ingles{order book}): a lista das ordens de compra e de venda que aguardam contraparte, organizada por preço. A Figura~\ref{fig:c11-livro} mostra um livro esquemático, com números inventados para o exemplo.

\figura[t]{cap11/livro_ofertas}{Livro de ofertas esquemático. Cada barra é a quantidade somada das ordens que aguardam em um preço. Não há ``o preço'' da ação: há o melhor preço para vender agora (20,00) e o melhor para comprar agora (20,02).}{fig:c11-livro}

A \termo{melhor oferta de compra} (\ingles{bid}) é o maior preço que alguém aceita pagar, e a \termo{melhor oferta de venda} (\ingles{ask}), o menor preço pelo qual alguém aceita vender. A distância entre as duas é o \termo{spread}, e o ponto médio entre elas é a referência usual de ``preço justo'' do instante. A quantidade disponível em cada nível é a \termo{profundidade}. As ordens são atendidas por prioridade de preço e, no mesmo preço, por ordem de chegada; por isso existe uma fila em cada nível \citep{harris2003}.

Quem tem pressa paga o spread. Uma compra imediata é executada contra o lado de venda, começando pelo melhor preço e avançando pelos seguintes até completar a quantidade. O Código~\ref{cod:c11-livro} faz essa conta para o livro da figura.

\lstinputlisting[style=py, firstline=8, lastline=17, caption={Uma ordem a mercado percorre o livro (\arq{cap11/l01\_livro.py})}, label=cod:c11-livro]{scripts/cap11/l01_livro.py}
\begin{lstlisting}[style=saida]
spread: 0.02 = 10.0 pb do preço médio 20.01
compra de  100: preço médio 20.0200,  5.0 pb acima do meio
compra de  500: preço médio 20.0200,  5.0 pb acima do meio
compra de 2000: preço médio 20.0290,  9.5 pb acima do meio
compra de 4000: preço médio 20.0358, 12.9 pb acima do meio
\end{lstlisting}

A unidade dessas contas é o \termo{ponto-base} (pb), um centésimo de ponto percentual: 10~pb são 0,10\%. Enquanto a ordem cabe no primeiro nível, o custo em relação ao preço médio é metade do spread, 5~pb. A compra de 2.000 ações esgota as 500 ofertadas a 20,02, as 1.200 a 20,03 e ainda leva 300 a 20,04; o preço médio sobe para 20,029 e o custo quase dobra. O spread mede o custo de uma ordem pequena e a profundidade determina o de uma ordem grande.

\subsection{Leilões de abertura e de fechamento}

O pregão começa e termina com um \termo{leilão}: durante alguns minutos as ordens se acumulam sem gerar negócios, e ao final as que puderem ser casadas são executadas a um preço único, aquele que maximiza a quantidade negociada. Os preços de abertura e de fechamento das ações líquidas saem desses leilões. Para quem simula com barras diárias, a consequência é que ``comprar na abertura'' ou ``vender no fechamento'' só é realizável participando do leilão, com a ordem enviada antes de o preço ser conhecido.

\section{Ordens: tipos, validade e ciclo de vida}

Uma ordem diz o que se quer (comprar ou vender, qual ativo, que quantidade) e em que condições. As condições de preço definem três tipos básicos, descritos em detalhe por \citet{harris2003}.

Uma \termo{ordem a mercado} manda executar já, pelo melhor preço disponível. Ela garante a execução e deixa o preço em aberto: é o caso do Código~\ref{cod:c11-livro}.

Uma \termo{ordem limitada} fixa o pior preço aceitável: uma compra limitada a 19,99 só é executada a 19,99 ou menos. Se não houver contraparte, ela entra no livro e espera. Garante o preço e deixa a execução em aberto. Quem usa ordem limitada deixa de pagar o spread e pode até recebê-lo, mas corre dois riscos. O primeiro é a não execução: o preço sobe sem voltar e a compra nunca acontece. O segundo é a \termo{seleção adversa}: a ordem parada no livro é uma oferta firme que qualquer um pode aceitar, e quem a aceita com mais frequência é quem sabe de algo. Suas compras limitadas tendem a ser executadas justamente quando o preço está caindo e vai continuar a cair.

Uma \termo{ordem stop} fica dormente até o mercado negociar a um preço de disparo. Um stop de venda a 9,50 é acionado quando sai um negócio a 9,50 ou menos, e nesse momento se transforma em outra ordem. No \emph{stop a mercado} ela vira ordem a mercado; no \termo{stop-limite} vira ordem limitada, e é preciso informar dois preços, o de disparo e o limite. O uso mais comum é limitar a perda de uma posição. Na B3, a oferta stop é do segundo tipo: ao ser disparada, entra no livro como oferta limitada \citep{b3ofertas}. O stop a mercado que aparece em algumas plataformas é montado pela corretora ou pela própria plataforma, por exemplo com um limite bem distante do disparo. A Tabela~\ref{tab:c11-ordens} compara os tipos.

\begin{table}[t]
\caption{O que cada tipo de ordem garante}
\label{tab:c11-ordens}
\begin{tabularx}{\textwidth}{>{\raggedright\arraybackslash}p{19mm} >{\raggedright\arraybackslash}p{25mm} >{\raggedright\arraybackslash}X >{\raggedright\arraybackslash}p{27mm}}
\toprule
\cab{Tipo} & \cab{Garante} & \cab{Deixa em aberto} & \cab{Uso típico} \\
\midrule
A mercado & execução imediata & o preço: paga o spread e, se for grande, percorre o livro & entrar ou sair com urgência \\
Limitada & o preço máximo de compra (mínimo de venda) & se e quando executa; seleção adversa & entrar com paciência; realizar lucro em um alvo \\
Stop a mercado & que a ordem é enviada quando o disparo ocorre & o preço, sobretudo em gap & limitar perda; entrar em rompimento \\
Stop-limite & disparo e preço-limite & a execução: em queda rápida pode não sair & limitar perda aceitando o risco de ficar na posição \\
\bottomrule
\end{tabularx}
\end{table}

\begin{armadilha}{Stop não garante preço}
O preço de disparo de um stop não é o preço de execução. Se a ação fecha a 10,00 e abre no dia seguinte a 8,00 depois de uma notícia, um stop de venda a 9,50 dispara na abertura e é executado perto de 8,00: nenhum negócio aconteceu entre 9,50 e 8,00, e o livro não tinha compradores ali. O stop-limite troca esse risco por outro, o de a ordem ficar parada acima do mercado enquanto o preço despenca. A Seção~\ref{sec:c11-simulador} mostra um caso real em que a diferença foi de 14\%.
\end{armadilha}

\subsection{Validade}

Toda ordem tem também um prazo. As duas validades que você encontra em praticamente qualquer plataforma são a \emph{do dia}, que expira no encerramento do pregão, e a \emph{até cancelar}. Bolsas e corretoras costumam oferecer ainda a validade até uma data e duas condições de execução imediata: \emph{executa ou cancela} (em inglês \ingles{immediate or cancel}, IOC), que executa o que for possível na hora e cancela o saldo, e \emph{tudo ou nada} (\ingles{fill or kill}, FOK), que é executada na hora pela quantidade inteira ou cancelada. O conjunto disponível muda conforme a corretora, a plataforma e o ativo; antes de programar, leia a documentação da sua.

\subsection{Ciclo de vida}

Depois de enviada, a ordem passa por estados, e o bot precisa acompanhar cada transição (Figura~\ref{fig:c11-estados}). A corretora e a bolsa primeiro validam a ordem: saldo, limites, preço dentro da faixa permitida. Se passar, ela é aceita e fica aberta. Daí pode ser executada de uma vez, executada aos pedaços, cancelada por você ou expirar.

\begin{figure}[t]
\centering
\begin{tikzpicture}[x=1mm, y=1mm, font=\sffamily\footnotesize,
  est/.style={draw, line width=0.4pt, minimum width=21mm, minimum height=8.5mm, align=center, inner sep=1pt},
  fim/.style={est, fill=lilas},
  seta/.style={-{Stealth[length=1.8mm]}, line width=0.4pt}]
  \node[est] (env) at (0,0) {Enviada};
  \node[est] (ace) at (30,0) {Aceita\\(aberta)};
  \node[est] (par) at (63,0) {Parcialmente\\executada};
  \node[fim] (exe) at (96,0) {Executada};
  \node[fim] (rej) at (0,-20) {Rejeitada};
  \node[fim, minimum width=27mm] (can) at (46.5,-20) {Cancelada\\ou expirada};
  \draw[seta] (env) -- (ace);
  \draw[seta] (ace) -- (par);
  \draw[seta] (par) -- (exe);
  \draw[seta] (env) -- (rej);
  \draw[seta] (ace.south) -- ([xshift=-7mm]can.north);
  \draw[seta] (par.south) -- node[right, text=cinza, pos=0.55] {\ saldo} ([xshift=7mm]can.north);
  \draw[seta] (ace.north) to[out=28, in=152] node[above, text=cinza] {de uma vez} (exe.north);
\end{tikzpicture}
\caption{Estados de uma ordem. Os estados finais estão sombreados. Uma ordem parcialmente executada continua aberta pelo saldo, que ainda pode ser executado, cancelado ou expirar.}
\label{fig:c11-estados}
\end{figure}

Dois detalhes desse diagrama causam a maior parte dos erros de quem começa. Enviar não é o mesmo que ter a ordem aceita: até a confirmação chegar, o programa não sabe se ela existe. E pedir o cancelamento não cancela na hora: o pedido viaja até a bolsa, e a ordem pode ser executada enquanto isso. O bot só deve alterar sua contabilidade quando recebe a confirmação de cada evento.

\section{Custos de transação}
\label{sec:c11-custos}

Cada operação tem custos de dois tipos. Os \termo{custos explícitos} aparecem na nota de corretagem:

\begin{itemize}
  \item \textbf{Corretagem}, cobrada pela corretora. Pode ser um valor fixo por ordem, um percentual do volume ou zero, conforme a corretora e o tipo de operação.
  \item \textbf{Emolumentos e taxas da bolsa}, cobrados pela B3 pela negociação e pela liquidação, como percentual do valor negociado. Para ações são da ordem de poucos pontos-base, com tabelas distintas para operações comuns e para as iniciadas e encerradas no mesmo dia (\ingles{day trade}).
  \item \textbf{Impostos} sobre o ganho líquido, com regras próprias de apuração, de compensação de prejuízos e de retenção na fonte.
\end{itemize}

Este livro não reproduz percentuais de propósito. As tarifas da B3 e as regras tributárias mudam, e um número impresso envelhece. Consulte a tabela de tarifas vigente no site da B3, a tabela de custos da sua corretora e a regra da Receita Federal em vigor, e coloque os valores atuais no seu modelo.

Os \termo{custos implícitos} não aparecem em nota nenhuma. Estão embutidos no preço de execução e, para quem gira a carteira com frequência, costumam ser a parte maior:

\begin{itemize}
  \item \textbf{Meio spread.} Comprar no \ingles{ask} e vender no \ingles{bid} custa metade do spread em cada ponta, medido contra o preço médio.
  \item \textbf{Slippage de latência.} O preço anda entre a decisão e a execução.
  \item \textbf{Impacto de mercado.} A própria ordem move o preço contra quem a envia: consome a profundidade e sinaliza ao mercado que há um comprador com pressa.
\end{itemize}

Um modelo de custo por operação soma as três partes que se conseguem estimar. Para uma ordem de valor $Q$, o custo por ponta (cada compra e cada venda é uma ponta), como fração do valor negociado, é
\begin{equation}
  c(Q) = \underbrace{c_{\text{taxas}}}_{\text{explícito}} + \underbrace{\frac{s}{2}}_{\text{meio spread}} + \underbrace{I(Q)}_{\text{impacto}} ,
  \label{eq:c11-custo}
\end{equation}
em que $s$ é o spread relativo e $I(Q)$ o custo de impacto, a única parcela que depende do tamanho. O custo padrão de 0,10\% por ponta dos outros capítulos, 10~pb, é um valor único para essa soma.

\subsection{A lei da raiz quadrada}

Para ordens pequenas em relação ao volume do dia, estudos com dados de execução de diferentes mercados encontram a mesma forma para o deslocamento do preço entre o primeiro e o último negócio de uma ordem executada aos poucos:
\begin{equation}
  \Delta(Q) \approx Y \, \sigma \sqrt{\frac{Q}{V}} ,
  \label{eq:c11-raiz}
\end{equation}
em que $\sigma$ é a volatilidade diária do ativo, $V$ o volume diário negociado e $Y$ uma constante de ordem um. A relação ficou conhecida como \termo{lei da raiz quadrada}. \citet{toth2011} a enunciam nessa forma e propõem uma explicação para ela; \citet{grinold2000} chegam à mesma forma por um argumento de risco de estoque de quem provê liquidez. Ela é uma regularidade empírica, e o expoente varia entre os estudos, de $0,4$ a $0,7$ segundo \citet{toth2011}; \citet{almgren2005}, com dados da mesa de ações de um grande banco americano, estimam $0,6$ para a parte temporária do impacto.

O deslocamento $\Delta$ é o impacto ao fim da ordem, e não o custo. Quem compra aos poucos paga, em média, menos que o preço final: se o preço sobe com a raiz quadrada da fração já executada, o preço médio de compra fica a dois terços de $\Delta$ acima do inicial. O custo de impacto da Equação~\eqref{eq:c11-custo} é, então,
\begin{equation}
  I(Q) \approx \tfrac{2}{3}\, Y \, \sigma \sqrt{\frac{Q}{V}} .
  \label{eq:c11-impacto}
\end{equation}
Use essas expressões para obter ordens de grandeza, e não previsões com casas decimais.

\begin{intuicao}
A forma côncava diz que os primeiros reais de uma ordem são os que mais movem o preço por real negociado. Quadruplicar a ordem dobra o impacto por ação. E a escala natural do tamanho é a fração do volume diário: R\$~1 milhão é uma ordem irrelevante em PETR4 e um evento em uma ação pouco negociada.
\end{intuicao}

O script \arq{cap11/l02\_custo\_ordem.py} aplica as Equações~\eqref{eq:c11-custo} e~\eqref{eq:c11-impacto} a PETR4 com os dados de 2025. A volatilidade e o volume saem dos arquivos; as taxas de 3~pb e a constante $Y = 1$ são hipóteses de trabalho, e o spread é suposto igual a um centavo, a menor variação de preço das ações na B3. A coluna \cod{fixo} soma taxas e meio spread; as três colunas de custo estão em pontos-base.

\begin{lstlisting}[style=saida]
preço 30.82  volume mediano R$ 983 mi  vol. diária 1.43%
   ordem (R$)  % do volume  fixo  impacto  total  custo (R$)
       10,000      0.001%   4.6      0.3    4.9           5
      100,000      0.010%   4.6      1.0    5.6          56
    1,000,000      0.102%   4.6      3.0    7.7         766
   10,000,000      1.017%   4.6      9.6   14.2      14,224
  100,000,000     10.170%   4.6     30.4   35.0     349,861
\end{lstlisting}

Para o investidor que negocia R\$~10 mil, o custo é quase todo fixo, cerca de 5~pb, e o impacto não existe na prática. Em R\$~10 milhões, 1\% do volume do dia, o impacto já é o dobro da parte fixa, e cada real negociado custa quase três vezes mais do que para a pessoa física. PETR4 é uma das ações mais negociadas da bolsa. Em ações menos líquidas a mesma conta, com $V$ dez ou cem vezes menor e spread de vários centavos, dá custos bem maiores para todos. O custo padrão de 10~pb do livro fica, por essa conta, do lado conservador para ordens pequenas em ações líquidas.

\section{Quanto custo uma estratégia aguenta}

O custo sai do retorno na proporção em que a estratégia negocia. O giro (\ingles{turnover}) dos Capítulos~\ref{cap:backtest} e~\ref{cap:estrategias} é o valor negociado em um período, somando compras e vendas, dividido pelo capital. Se cada real negociado custa $c$, o retorno líquido é
\begin{equation}
  R_{\text{líq}} = R_{\text{bruto}} - c \times \text{giro} .
  \label{eq:c11-liquido}
\end{equation}
O custo de equilíbrio (\ingles{break-even}) foi definido no Capítulo~\ref{cap:estrategias} como o custo por ponta que zera o ganho bruto acima da referência. Com $R^{f}$ a taxa da referência no período (o CDI),
\begin{equation}
  c^{*} = \frac{R_{\text{bruto}} - R^{f}}{\text{giro}} .
  \label{eq:c11-equilibrio}
\end{equation}
É essa a definição usada no livro. Em uma carteira comprada e vendida de soma zero, o retorno bruto já é um ganho acima da referência, e a expressão se reduz a $c^{*} = R_{\text{bruto}} / \text{giro}$. O número resume a folga da estratégia: se o custo que você estima pagar está acima de $c^{*}$, ela não se paga.

\subsection{Um experimento: o mesmo sinal, dois giros}

Para isolar o efeito do giro, o Código~\ref{cod:c11-custos} usa um único sinal em duas frequências de rebalanceamento. O sinal é o momentum 12--1 do Capítulo~\ref{cap:estrategias}, o retorno de 12 meses sem o mês mais recente, variante que se firmou depois de \citet{jegadeesh1993}, aplicado a seis ações (PETR4, VALE3, ITUB4, WEGE3, ABEV3 e BBAS3). A carteira compra as que subiram mais que a média e vende a descoberto as que subiram menos, com pesos proporcionais à distância da média, soma de pesos igual a zero e soma dos valores absolutos igual a um. Em uma versão os pesos são levados ao alvo todos os dias; na outra, só no último pregão de cada mês.

\lstinputlisting[style=py, firstline=10, lastline=28, caption={Momentum com rebalanceamento diário e mensal: retorno bruto e giro (\arq{cap11/l03\_custos.py})}, label=cod:c11-custos]{scripts/cap11/l03_custos.py}
\begin{lstlisting}[style=saida]
2011-01-11 a 2025-12-30
rebalanceamento         diário            mensal
giro anual (x capital)    24.8               5.3
custo por ponta    retorno  Sharpe   retorno  Sharpe
        0 pb         3.16%    0.19     2.88%    0.17
        5 pb         1.91%    0.11     2.62%    0.15
       10 pb         0.67%    0.04     2.35%    0.14
       20 pb        -1.81%   -0.11     1.83%    0.11
custo de equilíbrio    12.7 pb           54.5 pb
\end{lstlisting}

A função \cod{simular} recebe os pesos-alvo, os retornos e uma função que diz em que dias se negocia. Ela mantém o capital constante em 1 e acompanha o valor de cada posição: entre dois rebalanceamentos as posições derivam com os retornos, e o giro de um rebalanceamento é a soma dos valores absolutos das diferenças entre o alvo e o que a carteira tinha. A execução segue a convenção padrão do livro, no fechamento do sinal: os pesos calculados com o fechamento de um dia são assumidos nesse fechamento e rendem a partir do dia seguinte. Como a carteira tem soma zero, o que ela aplica nas compras vem das vendas, e seu retorno já é um retorno em excesso. Por isso o Sharpe aqui é a média sobre o desvio-padrão do próprio retorno, sem descontar o CDI, anualizado pela convenção de 252 dias (Seção~\ref{sec:anualizar}); a conta ignora o aluguel das ações vendidas a descoberto.

Sem custos, as duas versões rendem quase o mesmo, perto de 3\% ao ano. A diária negocia 24,8 vezes o capital por ano, quase cinco vezes mais que a mensal. Com o custo padrão de 10~pb por ponta, a diária perde quase quatro quintos do retorno e a mensal perde menos de um quinto. O custo de equilíbrio é de 12,7~pb em uma e 54,5~pb na outra, que é o que a Equação~\eqref{eq:c11-equilibrio} prevê: $3,16\% / 24,8 = 12,7$~pb. A Figura~\ref{fig:c11-sharpe} mostra as duas curvas.

\figura[t]{cap11/sharpe_custo}{Sharpe líquido do momentum de seis ações em função do custo por ponta. As duas versões partem de valores próximos; a inclinação de cada curva é aproximadamente proporcional ao giro. Os círculos vazados marcam o custo de equilíbrio.}{fig:c11-sharpe}

\begin{armadilha}{O Sharpe bruto deste exemplo não prova nada}
Um Sharpe de 0,19 em quase quinze anos tem erro-padrão próximo de 0,26 \citep{lo2002}: não se distingue de zero. Seis ações são muito pouco para uma estratégia de momentum. O experimento serve para mostrar a aritmética do custo, que vale para qualquer estratégia, e não para sugerir que esta dê dinheiro. A lição prática é outra: uma vantagem pequena só sobrevive com giro baixo, e antes de refinar o sinal convém perguntar quantas vezes por ano ele manda negociar.
\end{armadilha}

\subsection{Capacidade}

A Equação~\eqref{eq:c11-custo} torna o custo crescente com o tamanho, e daí sai um teto. A \termo{capacidade} de uma estratégia é o volume de capital que ela comporta antes de o impacto das próprias ordens consumir a vantagem. A Figura~\ref{fig:c11-impacto} junta as duas contas: a curva de custo de PETR4 da seção anterior e os custos de equilíbrio do experimento.

\figura[t]{cap11/impacto}{Custo por ponta em PETR4 em função do tamanho da ordem, sob as hipóteses da Seção~\ref{sec:c11-custos}. A versão diária do momentum deixa de se pagar com ordens acima de 0,7\% do volume do dia; a mensal suporta ordens dezenas de vezes maiores.}{fig:c11-impacto}

As últimas linhas da saída do mesmo script fazem a conta inversa: dado o custo de equilíbrio de cada versão, qual é a maior ordem que ainda se paga em um papel com a liquidez de PETR4.

\begin{lstlisting}[style=saida]
versão  equilíbrio  ordem máxima   % do volume  capital
diário     12.7 pb  R$    7.1 mi        0.72%  R$ 431 mi
mensal     54.5 pb  R$  270.1 mi       27.47%  R$ 3,803 mi
\end{lstlisting}

A versão diária deixa de se pagar com ordens acima de R\$~7 milhões em cada papel, 0,7\% do volume do dia. Para a mensal a conta dá R\$~270 milhões, 27\% do volume, um tamanho em que a lei da raiz quadrada, medida em ordens pequenas, já não merece confiança; leia o número como ``dezenas de vezes maior''. A última coluna converte ordem em capital, supondo seis papéis igualmente líquidos: a versão diária negocia perto de 10\% do capital a cada dia, e a mensal, 43\% a cada mês, de modo que a diferença de quase quarenta vezes no tamanho da ordem vira cerca de nove vezes em capital. São tetos otimistas, porque quatro das seis ações negociam menos que PETR4 e no teto o lucro já é zero: o ponto de lucro máximo em reais chega antes. A ordem de grandeza não muda com ajustes razoáveis das hipóteses: estratégias de giro alto são negócios pequenos.

\subsection{Executar uma ordem grande}

Quem precisa negociar uma fração relevante do volume do dia não envia uma ordem a mercado única, que percorreria o livro como no Código~\ref{cod:c11-livro}. Divide a ordem em fatias e as distribui no tempo. Os dois roteiros mais comuns têm nome. O \termo{TWAP} (\ingles{time-weighted average price}) usa fatias iguais em intervalos iguais. O \termo{VWAP} (\ingles{volume-weighted average price}) dimensiona as fatias na proporção do volume esperado em cada intervalo do dia, para que a ordem seja sempre a mesma fração do mercado; como o volume costuma se concentrar perto da abertura e do fechamento, um VWAP negocia mais nessas horas. A lei da raiz quadrada já descreve ordens executadas assim, aos poucos; fatiar não escapa dela, apenas evita o custo ainda maior de esvaziar o livro de uma vez.

Resta decidir em quanto tempo executar. \citet{almgren2001} formularam o dilema: executar depressa concentra as ordens e aumenta o impacto; executar devagar reduz o impacto, mas deixa a parte ainda não executada exposta à variação do preço durante mais tempo. Eles minimizam o custo esperado mais um múltiplo da variância do custo, com um parâmetro de aversão a risco que regula a troca. Sem aversão a risco, a solução é o ritmo constante do TWAP, que minimiza o impacto. Quanto maior a aversão a risco ou a volatilidade do ativo, mais a execução ótima se concentra no início do prazo, pagando impacto para reduzir a incerteza.

\section{Uma corretora simulada}
\label{sec:c11-simulador}

Antes de ligar um bot a uma corretora real, ele deve rodar contra uma corretora de mentira que responda como a real: aceite ordens, devolva estados, cobre custos e mantenha caixa e posição. A classe \cod{CorretoraPapel} do Código~\ref{cod:c11-papel} faz isso sobre barras diárias de abertura, máxima, mínima e fechamento (OHLC). Como uma barra diária esconde o que aconteceu dentro do dia, as regras de preenchimento precisam ser explícitas, e na dúvida escolhemos a hipótese pior para quem opera:

\begin{itemize}
  \item o bot manda processar a barra do dia antes de decidir, de modo que uma ordem só é executada em barras posteriores àquela em que foi criada;
  \item ordem \textbf{a mercado}: executa na abertura da barra seguinte, com um slippage contra você;
  \item ordem \textbf{limitada}: executa ao preço-limite apenas se a barra \emph{atravessar} o limite. Tocar não basta, nem na abertura, porque no preço exato haveria uma fila na sua frente. Se a barra abrir além do limite a seu favor, a ordem é tratada como ordem a mercado na abertura, com slippage, e o limite continua sendo o pior preço possível;
  \item ordem \textbf{stop}: é um stop a mercado. Dispara se a barra alcançar o gatilho e é executada ao gatilho com slippage; se a barra abrir além do gatilho, na abertura com slippage;
  \item barras sem volume ou sem preço não executam nada;
  \item toda execução paga uma comissão proporcional ao valor, e a corretora rejeita ordens malformadas, compras sem caixa e vendas sem posição.
\end{itemize}

\needspace{6\baselineskip}
\lstinputlisting[style=py, firstline=17, caption={Corretora simulada sobre barras diárias (\arq{cap11/papel.py})}, label=cod:c11-papel]{scripts/cap11/papel.py}

A classe \cod{Ordem}, definida no começo do arquivo, é um registro com identificador, lado, quantidade, tipo, preço e estado. O método \cod{\_preco} concentra as regras: a variável \cod{s} vale $+1$ para compra e $-1$ para venda, o que permite escrever cada regra uma vez para os dois lados. Executar na abertura supõe participar do leilão de abertura, e o slippage faz o papel do custo implícito da Seção~\ref{sec:c11-custos}; com ele, o preço de execução pode cair alguns pontos-base fora do intervalo da barra. As simplificações mais fortes são duas. Não há execução parcial: cada ordem é executada inteira ou não é. E as ordens não se conhecem: se o stop de uma posição é executado, o alvo continua aberto até que o bot o cancele.

\begin{repositorio}
O arquivo \arq{livro/scripts/cap11/teste\_papel.py} tem 14 testes, que cobrem as regras acima: limitada que apenas toca o preço, durante o dia ou na abertura, limitada que não sai mais barata que a ordem a mercado, stop com abertura em gap, stop e alvo na mesma barra, rejeições, barra sem volume. Rode-os antes de alterar a classe: uma corretora simulada com erro de contabilidade produz backtests errados sem dar aviso.
\end{repositorio}

\subsection{Um bot completo}

O Código~\ref{cod:c11-bot} liga um bot à corretora simulada, usando PETR4 de 2019 a 2024. A regra é simples de propósito e seus parâmetros foram fixados sem busca. Quando o fechamento está acima da média de 100 dias, o bot deixa uma compra limitada 2\% abaixo do fechamento, renovada todo dia. Executada a compra, envia duas ordens de saída: um stop 5\% abaixo do preço de entrada e uma venda limitada 8\% acima. Se nenhuma for executada em dez pregões, cancela as duas e sai a mercado. A quantidade é escolhida para que a perda planejada no stop seja 1\% do patrimônio, arredondada para lotes de 100 ações.

\lstinputlisting[style=py, firstline=1, lastline=38, caption={Bot de exemplo sobre a corretora simulada (\arq{cap11/l04\_bot.py})}, label=cod:c11-bot]{scripts/cap11/l04_bot.py}
\begin{lstlisting}[style=saida]
data       ordem lado   tipo      qtd   preço  taxa  pos.     caixa
2019-01-22 14    compra limitada 3000  6.5400  5.89  3000  80374.11
2019-02-07 17    venda  mercado  3000  6.6868  6.02     0 100428.46
2019-02-11 19    compra limitada 3100  6.4400  5.99  3100  80458.47
2019-02-14 21    venda  limitada 3100  6.9600  6.47     0 102028.00
2019-03-08 35    compra limitada 2900  6.8600  5.97  2900  82128.03
2019-03-15 37    venda  limitada 2900  7.4100  6.45     0 103610.58
2019-03-21 41    compra limitada 2700  7.4200  6.01  2700  83570.57
2019-03-25 42    venda  stop     2700  7.0081  5.68     0 102486.68
cancelada    529
executada    191
aceita         1
compra limitada: 96 executadas de 495 enviadas
venda stop: 35 executadas de 96 enviadas
venda limitada: 27 executadas de 96 enviadas
venda mercado: 33 executadas de 34 enviadas
pior stop: 2021-02-22, gatilho 7.19, executado a 6.1534
patrimônio final: 111,600.88   taxas pagas: 1,184.50
drawdown máximo: -5.6%   comprar e manter: +386%
\end{lstlisting}

O laço tem a forma da Figura~\ref{fig:c11-laco}: primeiro a corretora processa a barra do dia e executa o que estava pendente; depois o bot olha o fechamento e decide; por fim registra o patrimônio. As ordens criadas no passo 2 só podem ser executadas a partir da barra seguinte, e é essa ordem das chamadas que impede o bot de usar informação do futuro. Em relação à convenção padrão do livro, a execução no fechamento do sinal, o bot é mais conservador: decide com o fechamento do dia $t$ e só negocia a partir da abertura de $t+1$, como faz o backtrader no Capítulo~\ref{cap:backtest}. Fica entre a convenção padrão e a de um pregão inteiro de atraso usada no Capítulo~\ref{cap:projeto}.\nota{Os preços do arquivo são ajustados por proventos (Capítulo~\ref{cap:retornos}); por isso PETR4 aparece a R\$~6,54 em 2019. Preços e quantidades do extrato não são os de tela da época.}

No extrato, a primeira compra foi encerrada por tempo com uma ordem a mercado, a segunda e a terceira atingiram o alvo e a quarta foi estopada. Das 495 compras limitadas enviadas, 96 foram executadas, 11 delas na abertura, por um preço melhor que o limite; as outras foram canceladas e reenviadas a um novo preço no dia seguinte, o que em uma corretora real consome limite de mensagens.

A linha do pior stop merece atenção. Em 19 de fevereiro de 2021 o bot comprou a 7,57 e posicionou o stop a 7,19. No pregão seguinte a ação abriu 16,6\% abaixo do fechamento anterior, e o stop foi executado a 6,15, 14,4\% abaixo do gatilho. A operação perdeu 18,7\% em vez dos 5\% planejados, e o capital caiu 3,1\% em um dia, contra o risco de 1\% que o dimensionamento supunha. A Figura~\ref{fig:c11-bot} marca esse dia.

\figura[t]{cap11/bot}{Patrimônio do bot de exemplo, de 2019 a 2024. Os pontos marcam as saídas por stop, dimensionadas para custar 1\% do patrimônio; a de 22 de fevereiro de 2021 custou três vezes isso.}{fig:c11-bot}

O resultado em si é modesto: 11,6\% em seis anos, com o bot posicionado em 44\% dos dias e cerca de um quinto do patrimônio em cada operação, enquanto comprar e manter rendeu 386\% no mesmo período de alta forte. Como estratégia o bot não se recomenda. Ele serve de encanamento completo, que funciona igual para qualquer regra colocada no passo 2.

\subsection{O que uma barra não conta}

A barra diária registra quatro preços e nada sobre a ordem em que ocorreram. Para uma posição com stop e alvo, isso às vezes decide o resultado (Figura~\ref{fig:c11-barra}).

\figura[t]{cap11/barra_ambigua}{Uma barra, dois dias possíveis. A máxima passa do alvo e a mínima passa do stop. No caminho A a venda limitada é executada primeiro, com lucro; no caminho B o stop dispara antes, com prejuízo. O arquivo de barras é o mesmo nos dois casos.}{fig:c11-barra}

A \cod{CorretoraPapel} resolve a ambiguidade pela ordem em que as ordens foram enviadas, e o bot envia o stop antes do alvo: quando os dois poderiam ser executados na mesma barra, vale o stop, e a venda limitada é rejeitada por falta de posição. É a escolha conservadora. Na amostra do exemplo esse caso não ocorreu, mas outra limitação da mesma natureza sim: no dia da compra a posição ainda não tem stop, nem alvo, e em 2 das 96 entradas a mínima daquele mesmo dia ficou abaixo do nível em que o stop seria colocado (a máxima não passou do alvo em nenhuma). Quem precisa de respostas melhores usa barras de minutos ou os negócios um a um, o que muda o volume de dados, mas não a estrutura do simulador.

\subsection{Paper trading}

O passo seguinte é o \termo{paper trading}: o bot roda em tempo real, ligado aos dados ao vivo e a uma conta de simulação da corretora, sem dinheiro. Ele testa o que o simulador de barras não alcança: a conexão, o formato dos dados em tempo real, o tratamento de rejeições, o comportamento do programa na abertura e no fechamento, a rotina de ligar e desligar. Algumas semanas de paper trading sem intervenção manual costumam revelar a maior parte dos erros de programação.

\begin{armadilha}{Paper trading não mede custo}
A conta de simulação executa suas ordens contra o livro real sem alterá-lo. Suas ordens não consomem profundidade, não movem o preço e não ocupam lugar na fila; uma ordem limitada costuma ser dada como executada assim que o mercado toca o preço. A documentação da Alpaca, por exemplo, avisa que seu ambiente de simulação não reproduz impacto de mercado, slippage de latência nem posição na fila. O paper trading valida o programa. O custo só aparece com dinheiro real, e por isso se começa com posições pequenas.
\end{armadilha}

\section{Conectar-se a uma corretora}

As interfaces de programação (API) das corretoras variam nos detalhes e se parecem na estrutura:

\begin{itemize}
  \item \textbf{Autenticação por chave.} A corretora emite uma chave e um segredo, que identificam sua conta em cada requisição. Costuma haver chaves separadas para o ambiente de simulação e para o real.
  \item \textbf{REST para ordens e consultas.} Enviar, cancelar e consultar ordens, posição e saldo são requisições HTTP com resposta em JSON.
  \item \textbf{Websocket para tempo real.} Cotações, negócios e mudanças de estado das suas ordens chegam por uma conexão que fica aberta, sem que você precise perguntar.
  \item \textbf{Limites de requisição.} Há um teto de chamadas por minuto; quem o ultrapassa é bloqueado por um tempo. O bot precisa contar suas chamadas.
  \item \textbf{Identificador de ordem do cliente.} Você atribui a cada ordem um identificador único, e a corretora recusa uma segunda ordem com o mesmo identificador. É o que torna o envio \termo{idempotente}: se a resposta se perder e você reenviar, não compra duas vezes.
\end{itemize}

O bot não deve conhecer nenhum desses detalhes. Ele conversa com uma interface abstrata, e cada corretora ganha um adaptador que traduz a interface para a API correspondente. O Código~\ref{cod:c11-api} é um esboço que não foi executado: o endereço e os campos são fictícios, e uma versão real seguiria a documentação da corretora escolhida.

\needspace{6\baselineskip}
\begin{lstlisting}[style=py, caption={Esboço não executado de uma interface de corretora e de um adaptador REST}, label=cod:c11-api]
import os
from abc import ABC, abstractmethod

import requests


class Corretora(ABC):
    """Tudo o que o bot precisa saber pedir."""

    @abstractmethod
    def enviar(self, id_cliente, lado, qtd, tipo, preco=None): ...

    @abstractmethod
    def cancelar(self, id_cliente): ...

    @abstractmethod
    def estado(self, id_cliente): ...

    @abstractmethod
    def posicao(self, ativo): ...


class CorretoraREST(Corretora):
    URL = "https://api.corretora.exemplo/v1"     # fictício

    def __init__(self, ativo):
        chave = os.environ["CORRETORA_CHAVE"]    # nunca no código
        self.ativo = ativo
        self.http = requests.Session()
        self.http.headers["Authorization"] = f"Bearer {chave}"

    def enviar(self, id_cliente, lado, qtd, tipo, preco=None):
        corpo = {"id_cliente": id_cliente, "ativo": self.ativo,
                 "lado": lado, "qtd": qtd, "tipo": tipo,
                 "preco": preco}
        r = self.http.post(f"{self.URL}/ordens", json=corpo,
                           timeout=5)
        r.raise_for_status()                     # erro vira exceção
        return r.json()["estado"]
    # cancelar, estado e posicao seguem o mesmo padrão
\end{lstlisting}

O identificador \cod{id\_cliente} é criado pelo bot e gravado em disco \emph{antes} da primeira tentativa de envio. Se a requisição estourar o tempo, o bot reenvia com o mesmo identificador; gerar um novo a cada tentativa anularia a proteção. A \cod{CorretoraPapel} pode ser embrulhada na mesma interface, e então o mesmo bot roda no simulador, na conta de simulação e na conta real trocando uma linha.

\begin{armadilha}{Chave de API no repositório}
Uma chave de API com permissão de negociação equivale à senha da conta. Ela não entra no código nem no Git: fica em uma variável de ambiente ou em um arquivo listado no \arq{.gitignore}, como no Código~\ref{cod:c11-api}. Se uma chave for publicada por engano, apagar o arquivo não resolve, porque o histórico a guarda; revogue-a na corretora. Quando a corretora permitir, crie a chave sem permissão de saque e restrita ao seu endereço IP.
\end{armadilha}

\subsection{Caminhos usados por pessoas físicas}

Três caminhos aparecem com frequência, e os três têm documentação pública. Para o mercado brasileiro, algumas corretoras oferecem a plataforma MetaTrader~5, e a MetaQuotes, sua desenvolvedora, publica o pacote Python \cod{MetaTrader5}, que se comunica com o terminal instalado no computador (o pacote é distribuído para Windows) e expõe funções para obter cotações, consultar posições e enviar ordens. Para ações americanas, a Alpaca oferece uma API REST com ambiente de simulação gratuito. Para criptoativos, a biblioteca de código aberto \cod{ccxt} dá uma interface única para as APIs de mais de cem corretoras. Nomes de funções e parâmetros mudam entre versões: trabalhe com a documentação atual.


\section{Dimensionamento de posições}
\label{sec:c11-dimensionamento}

Falta a caixa \emph{Alvo de posição}: dado um sinal, quanto comprar. Três regras cobrem a maior parte dos usos.

A \termo{fração fixa} aplica em cada operação uma parcela constante do patrimônio, ou, na versão usada pelo bot do Código~\ref{cod:c11-bot}, arrisca uma parcela constante: a quantidade é o risco aceito dividido pela perda por ação até o stop. É simples e faz a posição encolher depois de perdas. Como se viu, o risco efetivo pode superar o planejado quando o stop é atravessado.

O alvo de volatilidade da Equação~\eqref{eq:c08-volalvo}, $w_t = \min(\sigma^{*}/\hat\sigma_t,\; w_{\max})$, ajusta a exposição para que o risco da posição fique constante. Com alvo de 10\% e um ativo que oscila 40\% ao ano, aplica-se um quarto do capital; se a volatilidade cai para 20\%, metade. A estimativa $\hat\sigma_t$ vem dos modelos dos Capítulos~\ref{cap:series} e~\ref{cap:risco}.

O \termo{critério de Kelly} responde a outra pergunta: qual fração maximiza o crescimento do capital no longo prazo? \citet{kelly1956} resolveu o caso de uma aposta repetida. Se a probabilidade de ganhar é $p$, a de perder é $q = 1 - p$ e cada real apostado rende $b$ reais em caso de ganho, a fração ótima é
\begin{equation}
  f^{*} = p - \frac{q}{b} .
  \label{eq:c11-kelly}
\end{equation}
Ela sai de maximizar o valor esperado do logaritmo da riqueza após uma aposta, $g(f) = p \ln(1 + b f) + q \ln(1 - f)$: derivando e igualando a zero chega-se à Equação~\eqref{eq:c11-kelly}. Com $p = 0,55$ e $b = 1$, $f^{*} = 0,55 - 0,45 = 0,10$.

Para um ativo com retorno em excesso esperado $\mu$ e volatilidade $\sigma$, a taxa de crescimento com alavancagem $f$ é aproximadamente $g(f) = f\mu - f^{2}\sigma^{2}/2$ acima da taxa livre de risco, e o máximo está em
\begin{equation}
  f^{*} = \frac{\mu}{\sigma^{2}} .
  \label{eq:c11-kelly-cont}
\end{equation}
Com $\mu = 6\%$ e $\sigma = 20\%$, $f^{*} = 0,06 / 0,04 = 1,5$: o critério manda alavancar uma vez e meia.

O script \arq{cap11/l05\_kelly.py} calcula $g(f)$ para a aposta binária do exemplo e simula 10 mil sequências de 500 apostas com três frações. Para cada uma imprime a taxa $g$, a riqueza final mediana, partindo de 1, e a mediana da maior queda sofrida no caminho:

\begin{lstlisting}[style=saida]
f* = 0.10
f = 0.05 (0.5 x Kelly): g = +0.00375  mediana final  6.53  queda 53%
f = 0.10 (1.0 x Kelly): g = +0.00501  mediana final 12.23  queda 83%
f = 0.20 (2.0 x Kelly): g = -0.00014  mediana final  0.93  queda 99%
\end{lstlisting}

A Figura~\ref{fig:c11-kelly} mostra a curva $g(f)$.

\figura[t]{cap11/kelly}{Crescimento esperado do logaritmo da riqueza em uma aposta com $p = 0,55$ e $b = 1$. A curva é chata perto do máximo e assimétrica nas consequências: apostar metade de Kelly custa um quarto do crescimento; apostar o dobro o elimina.}{fig:c11-kelly}

Na prática quase ninguém aposta o Kelly cheio, por duas razões que a saída deixa ver. A primeira é o caminho: no Kelly cheio, a trajetória mediana passa por uma queda de 83\%; com metade da fração a queda mediana é de 53\% e o crescimento esperado fica em três quartos do máximo ($0,00375$ contra $0,00501$). A segunda é o erro de estimação. As Equações~\eqref{eq:c11-kelly} e~\eqref{eq:c11-kelly-cont} supõem $p$ e $\mu$ conhecidos, e eles são estimados com muito ruído. Quem superestima a vantagem em duas vezes aposta o dobro de Kelly sem saber, e a última linha da saída mostra o destino: crescimento esperado nulo. Por isso se usa uma \termo{fração de Kelly}, tipicamente a metade ou menos; \citet{chan2021} discute o critério e a prática de operar com metade do valor de Kelly.

\section{Risco operacional}
\label{sec:c11-risco}

Um programa que envia ordens sozinho erra na velocidade em que opera. No fim de julho de 2012, a corretora americana Knight Capital instalou de forma incorreta uma versão nova do código de seu roteador de ordens, e uma função antiga e defeituosa, que nunca tinha sido removida, voltou a ser acionada. Segundo a \citet{sec2013}, nos primeiros 45 minutos do pregão de 1º de agosto o sistema enviou mais de 4 milhões de ordens ao tentar executar 212 ordens de clientes, negociou mais de 397 milhões de ações e deixou a empresa com um prejuízo superior a 460 milhões de dólares. O órgão regulador concluiu que faltavam controles capazes de barrar ordens erradas antes de chegarem ao mercado.

Um bot pessoal precisa das mesmas travas, em escala menor. Elas ficam na caixa \emph{Controle de risco} da Figura~\ref{fig:c11-laco} e são independentes do sinal.

\begin{itemize}
  \item \textbf{Verificação pré-ordem.} Antes de qualquer envio: a quantidade está abaixo do máximo por ordem? O preço está a poucos por cento do último negócio? A posição resultante respeita o limite? O número de ordens no último minuto é plausível? Uma resposta negativa bloqueia a ordem e gera um alerta.
  \item \textbf{Limites de posição e de perda diária.} Um teto de exposição por ativo e total, e uma perda no dia a partir da qual o bot para de abrir posições.
  \item \textbf{Botão de emergência} (\ingles{kill switch}). Um comando único que cancela todas as ordens abertas e impede novos envios, acionável por você e pelo próprio programa. Zerar as posições é uma decisão separada, que pode ou não ser automática.
  \item \textbf{Reconciliação.} A intervalos regulares, e sempre ao iniciar, o bot compara posição, caixa e ordens abertas com o que a corretora informa. Havendo diferença, vale a corretora, e o bot para até alguém entender a causa.
  \item \textbf{Registro} (\ingles{logging}). Cada dado recebido, decisão, ordem e resposta vai para um arquivo com horário. Sem registro não há como reconstituir o que aconteceu.
\end{itemize}

A queda de conexão merece um plano próprio. As ordens que já estão na bolsa continuam vivas enquanto o programa está fora do ar: um stop registrado na corretora protege a posição mesmo com o computador desligado, e um stop que existe só na memória do bot não protege nada. Ao reconectar, o programa não deve supor que nada mudou. Ele consulta ordens e posição, reconcilia e só então volta a decidir.

\subsection{Antes de ligar}

Operar com dinheiro real coloca você sob as regras da corretora e da Comissão de Valores Mobiliários (CVM). Leia o contrato e a política da corretora para ordens automatizadas, verifique se o uso da API é permitido para o seu tipo de conta e informe-se sobre a tributação das suas operações. Este livro é material de ensino: nada aqui é recomendação de investimento, e as estratégias mostradas servem para ilustrar métodos.

A lista abaixo reúne, na ordem em que devem ser cumpridas, as condições para um bot sair da simulação.

\begin{enumerate}
  \item O backtest passou pelos testes dos Capítulos~\ref{cap:backtest} e~\ref{cap:estrategias}, incluindo validação fora da amostra.
  \item O custo estimado por operação, com dados atuais de taxas e spread, está bem abaixo do custo de equilíbrio da estratégia.
  \item O bot rodou na corretora simulada e, depois, em paper trading por algumas semanas sem intervenção, com ordens iguais às que o backtest geraria nos mesmos dias.
  \item Limites de ordem, de posição e de perda diária estão configurados, e o botão de emergência foi testado.
  \item A reconciliação com a corretora roda ao iniciar e durante o dia; quedas de conexão e reinícios foram ensaiados.
  \item As chaves estão fora do código, sem permissão de saque.
  \item O capital inicial é pequeno, e há um critério escrito para aumentá-lo e outro para desligar o bot.
\end{enumerate}

No projeto final do Capítulo~\ref{cap:projeto}, os resultados devem ser apresentados líquidos de custos, com as ferramentas deste capítulo.

\begin{resumo}
  \item Um bot é um laço de eventos com responsabilidades separadas. Trabalhar com alvo de posição e reconciliar com a corretora permite reiniciar sem duplicar ordens.
  \item A ordem a mercado garante execução e não preço; a limitada garante preço e não execução, e sofre seleção adversa; o stop dispara uma ordem e não garante preço, sobretudo em gap.
  \item O custo por operação soma taxas, meio spread e impacto, e o impacto cresce aproximadamente com a raiz quadrada da fração do volume diário negociada.
  \item O retorno líquido é o bruto menos custo vezes giro. O custo de equilíbrio, ganho bruto acima da referência dividido pelo giro, diz se a estratégia tem folga; giro alto limita a capacidade.
  \item Simular com barras exige regras de preenchimento explícitas e conservadoras. O paper trading valida o programa, não os custos.
  \item Limites, verificação pré-ordem, botão de emergência, reconciliação e registro protegem contra o erro do próprio programa. Dimensione com alvo de volatilidade ou com uma fração do valor de Kelly.
\end{resumo}

\begin{exercicios}
  \item \nivel{básico} No livro da Figura~\ref{fig:c11-livro}, calcule à mão o preço médio de uma venda a mercado de 3.000 ações e o custo em pontos-base em relação ao preço médio de 20,01. Confira com a função \cod{varrer}.
  \item \nivel{básico} Uma estratégia rende, antes de custos, 8\% ao ano acima do CDI e gira o capital 40 vezes por ano, somando compras e vendas. Qual é seu custo de equilíbrio por ponta? Com custo de 12~pb, quanto sobra do ganho acima do CDI?
  \item \nivel{básico} Uma aposta paga R\$~2 por real apostado com probabilidade 0,4 e perde o valor apostado com probabilidade 0,6. Calcule a fração de Kelly e a taxa de crescimento $g$ em $f^{*}$, em $f^{*}/2$ e em $2 f^{*}$.
  \item \nivel{intermediário} Acrescente à \cod{CorretoraPapel} a validade ``do dia'': uma ordem não executada na primeira barra em que poderia ser é marcada como expirada. Escreva o teste antes de escrever o código.
  \item \nivel{intermediário} No Código~\ref{cod:c11-custos}, acrescente uma versão com rebalanceamento semanal e uma versão diária que só negocia quando a diferença entre o peso atual e o alvo passa de um limiar (uma \emph{banda de não negociação}). Como ficam o giro e o custo de equilíbrio?
  \item \nivel{intermediário} Rode o bot do Código~\ref{cod:c11-bot} com slippage de 0, 5, 20 e 50~pb e com a regra da ordem limitada trocada para ``executa se tocar''. Quanto do resultado depende de cada hipótese de preenchimento?
  \item \nivel{desafio} Embrulhe a \cod{CorretoraPapel} na interface \cod{Corretora} do Código~\ref{cod:c11-api} e acrescente ao bot um módulo de risco com limite de posição, limite de perda diária e botão de emergência. Escreva um teste em que um sinal defeituoso manda comprar todos os dias e mostre que o módulo o contém.
\end{exercicios}

\begin{leituras}
  \item \citet{harris2003}, capítulo 4, para os tipos de ordem e suas propriedades; o restante do livro cobre a microestrutura de mercado sem matemática pesada.
  \item \citet{chan2021}, capítulos 5 e 6, sobre sistemas de execução, paper trading e gestão de capital com o critério de Kelly.
  \item \citet{grinold2000}, capítulo 16, para a relação entre custos de transação, giro e valor adicionado de uma estratégia.
  \item \citet{almgren2001} e \citet{toth2011}, os artigos de referência sobre execução ótima e sobre a lei da raiz quadrada do impacto.
\end{leituras}
